\section{Workstream 3: data analysis}
\label{sec:analysis}

Report criterion 2 is worth 2 points, for analysis that is ``detailed and clear''. The brief
asks for analysis that will ``uncover patterns and insights'' and ``inform the selection of the
machine learning model''. So every analysis below is written as a question, then a plot, then a
statistic, then the decision it informs. The Week 5 visualisation notebook uses the same ``why
this plot / so what'' pattern.

\textbf{Where it lives.} \code{notebooks/01_eda.ipynb} calls \code{explorer.plots}, so every
figure is also produced by \code{python -m explorer.cli figures} as a vector PDF in
\code{reports/figures/}. The report includes the PDFs directly, never screenshots.

\textbf{Plot rules} (A1 NFR-01): a title; axis labels with units (BTC, satoshis, hours, count,
date UTC); the snapshot date in a footnote for DS1 plots; one colour per method across the whole
report; a colour-blind-safe palette; log axes wherever a variable spans orders of magnitude, with
that stated on the axis.

\subsection{DS1 and DS2 analyses}

\begin{longtable}{@{}L{3.6cm}L{3.4cm}L{3.0cm}L{4.4cm}@{}}
\toprule
\textbf{Question} & \textbf{Plot} & \textbf{Statistic} & \textbf{Decision it informs} \\
\midrule
\endhead
How big and how long is each history? & Bar chart of transactions per address; strip of history spans & median, IQR & Whether the 70/30 split leaves enough evaluation rows per address \\
What does ``normal'' look like for one address? & Daily timeline for 3 contrasting addresses (the Explorer's own chart) & -- & Motivates per-address baselines over global thresholds \\
How are F1--F4 distributed? & Histograms on log axes; box plots by direction & skewness, kurtosis & Heavy tails justify the log transforms, the robust z (median/MAD rather than mean/SD) and the spread floor \\
Do incoming and outgoing amounts differ? & Paired box plots of $|\text{net}|$ by direction per address & Mann--Whitney U per address & Evidence for CR-04 (same-direction size windows) \\
Are the features redundant? & Spearman correlation heatmap F1--F8 & $\rho$ & F1 and F4 are expected to correlate; justifies keeping both (absolute vs habitual size) or dropping one in models \\
Is activity periodic? & Hour-of-day and weekday histograms (UTC) & $\chi^2$ against uniform & Whether F8 (hour) is worth testing \\
How are gaps distributed? & Empirical CDF of gaps, log-x & quantiles & Sets the TxGPT gap-bucket edges; shows dormancy is rare \\
Does the feature space have dense groups and sparse points? & 2-D PCA of $(z_1..z_4)$, coloured by direction & explained variance & Non-spherical density supports DBSCAN alongside K-Means \\
How many transactions break $|z| \geq 3.5$ on the raw data? & Flag-rate bar per address & share per address & An early warning for the 20\% sanity rule; starts the threshold tuning \\
What kinds of addresses are in DS2? & Address-level feature pairs plot (transactions per day, in-share, median size, gap CV) & -- & Inputs to address-profile clustering; shows DS2 covers DS1's range, so TxGPT's training data matches its scoring data \\
How long are the DS2 sequences, and which tokens dominate? & Sequence-length histogram; token-frequency bars per field & -- & The context length (128 transactions) and the bucket design \\
\bottomrule
\end{longtable}

\subsection{DS3 and DS4 analyses}

\begin{longtable}{@{}L{3.6cm}L{3.4cm}L{3.0cm}L{4.4cm}@{}}
\toprule
\textbf{Question} & \textbf{Plot} & \textbf{Statistic} & \textbf{Decision it informs} \\
\midrule
\endhead
How imbalanced is Elliptic, and how does that change over time? & Stacked bars of illicit / licit / unknown per time step & illicit share per step & Temporal split; PR-AUC and F1 rather than accuracy; class weights \\
Which Elliptic features separate the classes? & Top-10 features by mutual information; violin plots & mutual information; Mann--Whitney U & LF versus AF comparison; RF importance cross-check \\
How imbalanced is BitcoinHeist, and in which years? & Positives per year and per family & counts & Fixes the temporal split year $Y$ and which families get a leave-out test \\
Are the BitcoinHeist features heavy-tailed? & Histograms on log axes by label & skewness & The $\log(1+x)$ transforms; robust scaling \\
Do ransomware and white rows differ? & Box plots of log income, length, weight and count by label & Mann--Whitney U with effect size (rank-biserial) & Which features carry signal; expectations for linear versus tree models \\
\bottomrule
\end{longtable}

\subsection{What the report's analysis section will contain}

Five findings, about 80 words each, each with one figure or a two-panel figure. The others stay
in the notebook and the appendix:
\begin{enumerate}
  \item Feature distributions are heavy-tailed, which motivates the log scales and robust z.
  \item Each address has its own scale, and in and out amounts differ, which motivates per-address
    and same-direction baselines.
  \item The z-feature space has dense cores and sparse points, which motivates DBSCAN alongside
    K-Means.
  \item The labelled data is severely imbalanced and drifts over time, which motivates temporal
    splits, class weights and PR-based metrics.
  \item DS2 covers the behaviour seen in DS1, which justifies training TxGPT on DS2 and scoring
    DS1.
\end{enumerate}
The wording above is a hypothesis. If the data says otherwise, the finding changes, and the
change itself is reported. That is still analysis that informed a decision.

\begin{done}
The notebook runs top to bottom from a clean kernel. Every figure has a ``so what'' sentence
under it. The five report figures exist as PDFs in \code{reports/figures/}. Statistical test
results are saved in \code{reports/tables/eda_tests.csv}.
\end{done}
